\documentclass[11pt]{article}

% ------------------------------------------------------------------
% Packages
% ------------------------------------------------------------------
\usepackage[utf8]{inputenc}
\usepackage[T1]{fontenc}
\usepackage{lmodern}
\usepackage[margin=1in]{geometry}
\usepackage{microtype}
\usepackage{setspace}
\usepackage{booktabs}
\usepackage{tabularx}
\usepackage{longtable}
\usepackage{array}
\usepackage{multirow}
\usepackage{enumitem}
\usepackage[table]{xcolor}
\usepackage{caption}
\usepackage{amsmath}
\usepackage{listings}
\usepackage[round,authoryear]{natbib}
\usepackage[colorlinks=true,linkcolor=blue!60!black,citecolor=blue!60!black,urlcolor=blue!60!black]{hyperref}
\usepackage{xurl}

\setstretch{1.12}
\setlist{itemsep=2pt,topsep=4pt}
\captionsetup{font=small,labelfont=bf}
\newcolumntype{L}{>{\raggedright\arraybackslash}X}
\newcolumntype{P}[1]{>{\raggedright\arraybackslash}p{#1}}
\renewcommand{\arraystretch}{1.18}

% Evidence labels
\newcommand{\fact}{\textcolor{green!45!black}{\textsf{\footnotesize[Fact]}}}
\newcommand{\likely}{\textcolor{orange!80!black}{\textsf{\footnotesize[Likely]}}}
\newcommand{\spec}{\textcolor{red!70!black}{\textsf{\footnotesize[Speculative]}}}
\newcommand{\verify}{\textcolor{gray}{\textsf{\footnotesize[verify]}}}

\lstset{basicstyle=\ttfamily\footnotesize,columns=fullflexible,breaklines=true,frame=single,
        framerule=0.3pt,rulecolor=\color{gray!60},backgroundcolor=\color{gray!5},
        commentstyle=\color{gray}}

% ------------------------------------------------------------------
\title{\textbf{Who Are Colombia's NINIs?}\\[4pt]
\large Measuring and Explaining Youth Disconnection from School and Work\\[2pt]
\normalsize Proof of Concept}
\author{Juliana S\'anchez}
\date{November 2026 \quad--\quad Working draft, not for circulation}

\begin{document}
\maketitle

\begin{abstract}
\noindent About one in five Colombians aged 15--28 neither studies nor works. Two thirds of them are women, even though women now outnumber men in tertiary education. Colombia also reports one of the highest NEET rates in the OECD. Yet the label ``NINI'' bundles very different situations: young mothers doing full-time unpaid care, recent high-school graduates waiting to enrol, unemployed job seekers, people with disabilities, and a residual group that is genuinely disconnected. This note (i) summarises what is known, (ii) proposes a measurement framework that decomposes the NINI stock into economically meaningful groups and separates transitory from persistent spells, (iii) lists the open questions and puzzles, (iv) states competing explanations of \emph{why} young people become NINIs, with a prediction that distinguishes each, (v) catalogues vocational-education and other policy changes that offer variation to test them, (vi) reviews the academic, official and international literature, including a full review of the central bank's synthetic-cohort study, and (vii) maps the available data and a first empirical step. The main finding of the literature review is that the Colombian evidence is thin and almost entirely descriptive or correlational: no study identifies causal determinants or consequences of NINI status, and no Colombian experiment uses NINI status as an outcome.
\end{abstract}

\noindent\fbox{\parbox{\dimexpr\linewidth-2\fboxsep-2\fboxrule}{\small
\textbf{How to read this note.} Claims are tagged \fact\ (documented in a source), \likely\ (supported but involving inference) or \spec\ (plausible, untested). The literature and figures were assembled through web search in November 2026; most PDFs could not be opened from the research environment, so figures come from source summaries and press reports that cite the institutions. Items tagged \verify\ or with a ``Verify'' note in the bibliography must be checked against the original before citing.}}

\tableofcontents
\newpage

% ==================================================================
\section{Motivation}
% ==================================================================

Three facts motivate the project.

\begin{enumerate}
\item \textbf{Size.} DANE's monthly youth labour-market bulletin reports roughly 2.3--2.7 million NINIs aged 15--28 in 2025--2026, i.e.\ 21--24\% of the age group \fact\ \citep{dane2026}. The OECD places Colombia second among member countries for the 18--24 NEET rate (27.1\% in \textit{Education at a Glance 2022}, behind Turkey) \fact\ \citep{oecd2022}.
\item \textbf{Gender paradox.} About two NINIs in three are women in every source and definition \fact\ \citep{dane2026,ilo2024col,oecd2022}, while women dominate access to and graduation from tertiary education. Young women are 2.1 times more likely than young men to be NINI \citep{oecdgender}. How can women be simultaneously the most educated and the most disconnected group?
\item \textbf{Thin evidence.} The Colombian literature consists of about six microdata studies, all cross-sectional logit or probit models on household surveys, plus one pseudo-panel by the central bank \citep{lasso2026}. None identifies causal mechanisms (Section~\ref{sec:lit}).
\end{enumerate}

The central argument of this note is that \emph{measurement comes first}. ``NINI'' is a policy label, not an economic state. Before asking why the rate is high or why it is female, we need to know how much of it is unpaid care (which is work), how much is job search, how much is a short transition after school, and how much is persistent disconnection. Several apparent puzzles may be composition effects \likely.

% ==================================================================
\section{Background}
% ==================================================================

\subsection{Stylised facts}

Table~\ref{tab:facts} collects headline figures. They are not directly comparable: they differ in age range, definition, coverage and period (Section~\ref{sec:measure}).

\begin{table}[htbp]
\centering\small
\caption{Headline NINI figures for Colombia}\label{tab:facts}
\begin{tabularx}{\linewidth}{P{2.4cm}P{3.2cm}P{3.0cm}L}
\toprule
Period & Source & Age / definition & Figure \\
\midrule
2019 & DANE (via press) & 15--28, broad & $\approx$22\% \\
May--Jul 2020 & DANE (via press) & 15--28, broad & 33\% (COVID peak) \\
2020 & DANE (via press) & 15--28, broad, within sex & 42\% of women vs.\ 23\% of men \\
2008--2021 avg. & \citet{lee60} & 14--28, GEIH & $\approx$26\%; gender gap 18.6 pp \\
2015--2018 & \citet{lopez2020} & youth, household surveys & 25\% of women vs.\ 8\% of men \\
2022 & \citet{oecd2022} & 18--24, OECD & 27.1\%; 67\% women \\
2023 & \citet{lee99} & 15--28, \emph{strict} (not seeking work) & 13.8\% (DANE broad: 22\%) \\
Q1 2024 & \citet{ode2024} & 15--28, Bogot\'a & 18.5\%; 64\% women; women do $\approx$30 h/week unpaid work \\
1984--2024 & \citet{lasso2026} & 15--24, seven cities & fell $\approx$7 pp to 17--18\% in 2024 (the paper reports both 18.2\% and 17.2\%); $\approx$5\% in top- vs.\ $\approx$40\% in bottom-quintile households \\
2016--2023 & \citet{ilo2024col} & 15--28 & women/men ratio fell from 2.4 to 1.9 \\
Dec 2025--Feb 2026 & \citet{dane2026} series & 15--28, broad & 24.2\% (2.67 M); 66\% women \\
Jun--Aug 2026 & \citet{dane2026} & 15--28, broad & $\approx$21\% (2.33 M of 11.07 M) \\
\bottomrule
\end{tabularx}
\par\smallskip\raggedright\footnotesize Notes: ``Broad'' = not studying and not employed (includes the unemployed); ``strict'' = also not seeking work. Press-reported DANE figures should be checked against the bulletins. Rolling quarters are seasonal (December--February is the highest), so only like-for-like quarters should be compared.
\end{table}

Four regularities emerge \likely: (i) roughly two women per man; (ii) a steep income gradient; (iii) a sharp jump in the rate between ages 16 and 18, after which the cohort NINI \emph{share} remains above 15\% \citep{lasso2026}, which does not by itself show that the same individuals remain NINI; and (iv) a territorial gradient, from about 18\% in Bogot\'a to 27--29\% in Cartagena, with Valledupar and Riohacha reported as the highest cities.

\subsection{Institutional context}

\begin{itemize}
\item \textbf{Legal definition of youth.} Law 1622 of 2013 defines youth as ages 14--28. DANE publishes NINI figures for 15--28, following the ILO working-age threshold. The SDG indicator 8.6.1 uses 15--24 \citep{sdg861}; the OECD often reports 18--24.
\item \textbf{Survey redesign.} The GEIH was redesigned with the 2018 census frame. In October 2022 DANE released back-cast (\textit{empalmadas}) series; pre-2022 figures are only comparable through those series \fact.
\item \textbf{Policy framework.} CONPES 173 (2014) and CONPES 4040 (2021, \textit{Pacto Colombia con las Juventudes}) rely on DANE diagnostics \citep{conpes173,conpes4040}. The 2018--2022 National Development Plan states that 71\% of NINIs are women. ICBF produced a characterisation of NINIs aged 14--28 with three profiles \citep{icbf}.
\item \textbf{Programmes targeting this population.} J\'ovenes en Acci\'on (randomised evaluations exist; Section~\ref{sec:lit}), Sac\'udete (state policy since Law 2231 of 2022, operated by ICBF), J\'ovenes en Paz (Decree 1649 of 2023) and the free-tuition policy for public higher education. To our knowledge none has been evaluated with NINI status as an outcome.
\end{itemize}

% ==================================================================
\section{Defining and measuring NINIs}\label{sec:measure}
% ==================================================================

\subsection{The conceptual problem}

The NEET indicator has no international standard and is easily misinterpreted \citep{elder2015}. The SDG metadata itself warns that the prevalence and composition of each subgroup have different policy implications \citep{sdg861}. In Colombia, the indicator mixes at least two very different populations: inactive young women doing care work, and unemployed young men \citep{ochoa2015,hernandez2016}. A conference abstract reports that 88.5\% of female NEETs are inactive versus 30\% of male NEETs in Latin America \verify. A regional survey of 15,000 young people in nine countries, Colombia included, finds that only about 3\% do no productive activity at all (no study, work, care or job search) and have no disability \citep{novella2018}.

\subsection{Existing definitions}

\begin{table}[htbp]
\centering\small
\caption{Definitions in use and why figures differ}\label{tab:defs}
\begin{tabularx}{\linewidth}{P{3.3cm}LP{4.2cm}}
\toprule
Dimension & Variants & Who uses it \\
\midrule
Age range & 14--28 (Law 1622); 15--28; 15--24; 18--24 & ICBF, LEE 60; DANE; SDG/ILO, BanRep, IDB; OECD \\
Labour status & Broad (not employed, incl.\ unemployed) vs.\ strict (inactive only, ``NLEET'') & DANE, ILO, OECD vs.\ LEE 99, LaboUR \\
Education & Formal enrolment only vs.\ also SENA / non-formal training & Varies; rarely documented \\
Coverage & National; 13 or 7 cities; single city & DANE; BanRep; local observatories \\
Period & Rolling quarter vs.\ annual average & DANE bulletins vs.\ academic papers \\
Population base & Pre-2022 frame vs.\ 2018-census frame & Use back-cast series \\
\bottomrule
\end{tabularx}
\end{table}

The gap between the strict 13.8\% \citep{lee99} and the broad 22\% for 2023 (about 8 pp) is almost entirely unemployed job seekers \likely. The female share is higher under the strict definition (75\% or more) than under the broad one (64--68\%), consistent with women's NINI status being driven by inactivity and men's by unemployment \likely.

\subsection{Proposed typology}

Instead of a single rate, we propose to decompose the stock into mutually exclusive groups, built sequentially from the GEIH and validated with the time-use survey (ENUT):

\begin{table}[htbp]
\centering\small
\caption{Proposed NINI typology}\label{tab:typology}
\begin{tabularx}{\linewidth}{P{3.2cm}LP{4.3cm}}
\toprule
Group & Definition & Source of classification \\
\midrule
\textbf{N1. Searching} & Not studying, not employed, actively seeking work & Labour-force status (unemployed) \\
\textbf{N2. Transition} & Inactive, aged $\le$19, completed secondary within the last year or awaiting enrolment & Age, completed level, reason for inactivity \\
\textbf{N3. Care} & Inactive; main activity is household chores/care; split by own child $<$5 and partnership & Main activity, household roster \\
\textbf{N4. Disability / health} & Inactive due to permanent disability or illness & Main activity, reason for not seeking \\
\textbf{N5. Discouraged} & Inactive, wants to work, not searching because no jobs are available & Reason for not seeking \\
\textbf{N6. Disconnected} & Residual: none of the above & Residual \\
\bottomrule
\end{tabularx}
\end{table}

The care group should be validated with ENUT data on hours of unpaid care (2012--13, 2016--17, 2020--21). A natural benchmark is the ``truly idle'' share: N6 as a fraction of all youth, which may be far below the headline 21\% \spec.

\subsection{Stocks versus flows}

The GEIH is a repeated cross-section. A stable 21\% rate is consistent both with a small group that is permanently disconnected and with many young people passing briefly through the state. Distinguishing the two requires longitudinal information:

\begin{itemize}
\item \textbf{Synthetic cohorts.} Following birth cohorts across survey years, as in \citet{lasso2026}.
\item \textbf{Household panel.} The Colombian Longitudinal Survey (ELCA, Universidad de los Andes) follows households over several rounds \verify\ (confirm the rounds available).
\item \textbf{Administrative linkage.} Saber 11 test-takers can be followed into tertiary enrolment (SPADIES/SNIES), SENA and formal employment (PILA). Never appearing in enrolment or PILA is an \emph{upper bound} on NINI status, because informal workers are invisible in PILA.
\end{itemize}

\subsection{Measurement pitfalls}

\begin{itemize}
\item \textbf{Employment threshold.} One hour of paid work in the reference week counts as employed; unpaid family work counts as employment.
\item \textbf{What counts as studying.} Pre-university courses, SENA and non-formal training may or may not be recorded as attendance.
\item \textbf{Excluded populations.} The GEIH samples private households. Young men doing compulsory military service live in barracks and are presumably outside the sample \spec; this could mechanically affect male rates at ages 18--20.
\item \textbf{Denominators.} Population projections based on the 2018 census, which had high omission, enter the expansion factors.
\item \textbf{Seasonality.} Rolling quarters around the end of the school calendar inflate the rate.
\end{itemize}

% ==================================================================
\section{Research questions and puzzles}\label{sec:questions}
% ==================================================================

\subsection{What we know}

\begin{itemize}
\item Roughly two in three NINIs are women, in every source \fact.
\item Female NINIs are associated with early partnership, children under five and domestic work; male NINIs with job search \fact\ (correlational) \citep{ochoa2015,lopez2020,lenis2021}.
\item The rate jumps from about 5\% to 22\% between ages 16 and 18; afterwards the cohort NINI share stays high. Cohorts with higher NINI shares at 16--22 show lower employment (men) and less salaried work (women) at 36--42 \fact\ (cohort-level, correlational) \citep{lasso2026}.
\item The NINI rate is countercyclical, peaking in 1999 and 2020 \fact\ \citep{lasso2026}.
\item The rate is about 5\% in high-income and 40\% in low-income households \fact\ \citep{lasso2026}.
\end{itemize}

\subsection{What we do not know}

Table~\ref{tab:puzzles} lists the questions, ordered by how informative a first descriptive step would be. Hypotheses and designs are inferences, not findings.

\begingroup\small
\begin{longtable}{P{3.3cm}P{4.4cm}P{4.1cm}P{2.8cm}}
\caption{Open questions, candidate hypotheses and sources of variation}\label{tab:puzzles}\\
\toprule
Question & Candidate hypotheses & Variation that could discriminate & Data \\
\midrule
\endfirsthead
\toprule
Question & Candidate hypotheses & Variation that could discriminate & Data \\
\midrule
\endhead
\bottomrule
\endfoot
\textbf{Q1. Is the gender paradox composition or behaviour?} & (a) Different populations: university women are mostly middle class; NINI women are poor, partnered early, mothers. (b) A care penalty that also affects educated women. (c) Norms; (d) weak demand for low-skilled female labour; (e) lack of childcare. & Decomposition by income, education and motherhood; within-group comparisons of women with and without children; staggered childcare expansion (CDI rollout). & GEIH, ENUT, ENDS \\
\textbf{Q2. Why the jump at ages 16--18?} & Colombian secondary school ends after grade 11, so many graduate at 16--17, before 18, the age many formal jobs require \spec. Tertiary immediate transition is only about 40\%. & Age in months around graduation; date-of-birth vs.\ school-entry cut-offs; comparison of graduates and dropouts at the same age. & GEIH (age, level), Saber 11 cohorts \\
\textbf{Q3. How much is transitory vs.\ persistent?} & A large share of spells may be short post-graduation gaps; a smaller core is chronic. & Cohort tracking; duration of search; prior work experience. & Saber 11--SPADIES--PILA linkage; ELCA \\
\textbf{Q4. Why did the women/men ratio fall from 2.4 to 1.9 (2016--2023)?} & Women exiting faster (ILO) vs.\ men entering; falling fertility; post-COVID male job losses. & Decompose changes in the female and male rates separately by age and motherhood. & GEIH 2016--2025 \\
\textbf{Q5. Is Colombia's high OECD ranking real?} & Definitions and coverage (conscripts, rural inactivity, unpaid family work classification) vs.\ genuine disconnection. & Re-compute OECD/SDG definitions from GEIH microdata; compare with SEDLAC harmonised series. & GEIH, SEDLAC \\
\textbf{Q6. What are the consequences?} & Scarring in formal employment and wages; earlier fertility; crime; mental health. & Local labour-demand shocks at the time of leaving school (Gregg-style instruments); cohort exposure to COVID. & PILA, birth records, police data \\
\textbf{Q7. Why the territorial gradient?} & Local labour demand, conflict and displacement, ethnicity, early unions in the Caribbean. & Municipal variation in the 2018 census; shift-share labour-demand shocks. & Census 2018, GEIH \\
\textbf{Q8. Did COVID leave a cohort scar?} & The 2020 graduating cohort missed a transition window. & Compare adjacent graduating cohorts over time. & GEIH, PILA \\
\textbf{Q9. What about Venezuelan migrant youth?} & Legal status, recognition of credentials, care burdens. & Regularisation timing; municipal arrival intensity. & GEIH migration module \\
\textbf{Q10. Do current policies reduce NINI status?} & Free tuition, Renta Joven, J\'ovenes en Paz, Sac\'udete. & SISBEN eligibility cut-offs (RD); staggered municipal rollout. & SISBEN, SIMAT, PILA, GEIH \\
\end{longtable}
\endgroup

% ==================================================================
\section{Why do young people become NINIs? Competing explanations}\label{sec:why}
% ==================================================================

The existing evidence describes \emph{who} and \emph{when}, but not \emph{why}. Explanations in the Colombian literature are either restatements of the definition (``low tertiary coverage and low labour insertion''), mechanisms borrowed from other countries, or correlations with aggregate variables (Section~\ref{sec:lasso}). The research question we propose is therefore explicitly causal: \emph{which mechanisms move young Colombians into, and keep them in, NINI status, and do they differ by sex?}

The strategy is to state competing explanations and, for each, a prediction that the others do not share. Table~\ref{tab:why} is the core of the project. All entries are hypotheses \spec\ until tested.

\begingroup\small
\begin{longtable}{P{2.8cm}P{5.0cm}P{4.4cm}P{2.6cm}}
\caption{Competing explanations and discriminating predictions}\label{tab:why}\\
\toprule
Explanation & Prediction that sets it apart & Variation / test & Data \\
\midrule
\endfirsthead
\toprule
Explanation & Prediction that sets it apart & Variation / test & Data \\
\midrule
\endhead
\bottomrule
\endfoot
\textbf{E1. Liquidity constraints} (the central bank paper's main interpretation) & NINI rises when household income falls around graduation; falls when tuition becomes free or youth transfers start; effect concentrated in low-SES youth who graduated (not dropouts). & Free tuition in public higher education; SISBEN cut-offs of youth transfers (RD); household income shocks at 16--18. & GEIH, SPADIES, SISBEN \\
\textbf{E2. Low perceived returns or misinformation} & NINI responds to local returns to tertiary education and to information, not to income. & City-level college premium over time; distance to tertiary supply; information interventions. & GEIH, SNIES \\
\textbf{E3. Weak youth labour demand and legal-age barriers} & The 16--18 jump is concentrated among graduates who are still under 18; NINI responds to local youth labour demand and to the bite of the minimum wage. & Age at graduation (date of birth vs.\ school-entry cut-off); Bartik demand shocks by city; minimum-wage bite (Kaitz index) by city. & Saber 11, GEIH \\
\textbf{E4. Care, motherhood and early partnership} & Female NINI is concentrated in ``household chores''; rises sharply after first birth or union; falls when childcare supply expands. & Event study around first birth; staggered expansion of early-childhood centres (CDI); ENUT care hours. & GEIH, ENUT, ENDS, vital statistics \\
\textbf{E5. School quality and dropout} & NINI concentrated among those who leave before completing secondary; responds to school-level shocks (teacher supply, school day). & Completed level; school-level shocks (full school day, teacher hiring). & SIMAT, GEIH \\
\textbf{E6. Training supply and fit} & NINI falls where and when vocational training supply (SENA places, articulation with upper secondary) expands. & Staggered SENA and articulation rollout across municipalities and schools (Section~\ref{sec:policy}). & SENA, SIMAT, GEIH \\
\textbf{E7. Discouragement or preferences} & Inactive youth report wanting to work but not searching; no response to income, demand or supply shocks. & Reason for not searching; heterogeneity across shocks. & GEIH \\
\textbf{E8. Measurement} & Gaps shrink under the strict definition, when care is counted, or when coverage is national. & Recompute under each definition (Table~\ref{tab:defs}). & GEIH, ENUT \\
\end{longtable}
\endgroup

\paragraph{First step.} Before any causal design, the descriptive decomposition (Section~\ref{sec:plan}) plus an individual-level model of NINI status by type---which the central bank paper does not estimate---already rules some explanations in or out. For example, if most female NINIs are in the care group and the jump at 16--18 is concentrated among graduates under 18, E4 and E3 deserve causal designs before E1.

% ==================================================================
\section{Policy changes and vocational education: candidate sources of variation}\label{sec:policy}
% ==================================================================

Colombia changed many rules that sit exactly on the school-to-work margin. Table~\ref{tab:policy} lists the most relevant ones, the explanation in Table~\ref{tab:why} they speak to, and the variation they offer. Two findings from the policy review stand out. First, \textbf{no Colombian evaluation of any of these policies uses NINI status as an outcome} \likely\ (none found in English or Spanish searches). Second, the OECD reports an average graduation age from vocational upper secondary in Colombia of 16, against 22 in the OECD \fact\ \citep{oecd2022vet}: the 16--18 NINI jump coincides with leaving upper secondary, so policies acting in grades 10--11 or just after graduation are on the relevant margin.

All facts in the table come from search snippets of official and press sources (November 2026) and should be checked against the legal texts; items marked \verify\ are uncertain.

\clearpage
\begingroup\small\setlength{\tabcolsep}{3pt}
\begin{longtable}{P{2.6cm}P{4.0cm}P{0.9cm}P{3.4cm}P{3.1cm}P{1.2cm}}
\caption{Policy changes affecting the school-to-work transition}\label{tab:policy}\\
\toprule
Policy & Date and rule & Expl. & Exploitable variation & Existing evidence & Prom. \\
\midrule
\endfirsthead
\toprule
Policy & Date and rule & Expl. & Exploitable variation & Existing evidence & Prom. \\
\midrule
\endhead
\bottomrule
\endfoot
\multicolumn{6}{l}{\textit{Vocational education and training}}\\
J\'ovenes en Acci\'on redesign; Renta Joven & DPS Res.\ 01970 (Nov 2012): transfer conditional on SENA or tertiary enrolment for poor secondary graduates, SISBEN-based. Replaced by Renta Joven on 1 Jan 2024 (ages 14--28, secondary completed). & E1, E6 & SISBEN eligibility cut-off (RD); municipal phase-in \verify & RD: eligibility raises SENA enrolment by 8--13 pp \citep{gamboa2021}. NINI outcome not studied. & High \\
Articulation of upper secondary with SENA (double certification, grades 10--11); Bogot\'a \textit{Media Fortalecida} & School-by-school adoption from the early 2000s \verify; Bogot\'a Project 891 (2012--2016), $\approx$85,000 students by 2015 and an optional grade 12 in 21 schools. & E6, E3 & Staggered adoption by school (staggered DiD on graduates' later enrolment, work and NINI status). & DNP evaluation compares participants and non-participants (method unknown); no causal study found. & High \\
Apprenticeship contract & Law 789 of 2002: firms with $\geq$15 workers must hire apprentices; stipend 75\% of minimum wage, 100\% if prior-year unemployment $<$10\%. Law 2466 of 2025: apprenticeship becomes a labour contract (75\%/100\%), monetisation at 1.5 minimum wages per missing apprentice. & E3, E6 & Firm-size threshold; unemployment-triggered stipend rule; 2025 reform date. & RD at the 15-worker threshold for firm outcomes \citep{ospino2018}. SENA reports 392,056 apprentices in 2024 and 379,800 in 2025 (disputed). & Med.--High \\
SENA expansion & Instructors rose from 6,166 to 18,688 (2003--2010); training hours more than doubled (2002--2005); large virtual-course expansion. & E6 & Municipal variation in SENA supply over time \verify\ (series not located). & No causal study found. & Medium \\
\newpage
\multicolumn{6}{l}{\textit{Schooling}}\\
Repeal of the 5\% repetition cap & Decree 230 of 2002 required $\geq$95\% promotion; Decree 1290 of 2009 removed the cap from the 2010 school year. & E5 & School-level change in retention rates (DiD). & Retention raised same-year dropout, most at grades 9 and 11 \citep{ferreira2020}. Destination of dropouts (work vs.\ NINI) unknown. & High \\
Free public schooling & Decree 4807 of 2011: free grades 0--11 nationwide from 2012; earlier local schemes (Bogot\'a 2004 for the poorest SISBEN level). & E1, E5 & Pre-existing local free schooling as dose; Bogot\'a SISBEN cut-off. & RD for Bogot\'a \citep{barrera2007}; national reform not evaluated. & Medium \\
Full school day (\textit{Jornada \'Unica}) & Law 1753 of 2015 art.\ 57; Decrees 501 of 2016 and 2105 of 2017; phased in by school, campus and zone. & E5 & Staggered adoption by school. & Earlier full-day schooling raised test scores $\approx$0.1 SD \citep{hincapie2016}; nothing on dropout or NINI. & Medium \\
Tertiary financing & Ser Pilo Paga (2014--2018: Saber 11 and SISBEN cut-offs); Generaci\'on E (2019); Matr\'icula Cero (2020--21); Law 2307 of 2023 (free public undergraduate tuition). & E1 & Score and SISBEN cut-offs (RD); timing of free tuition by university type. & Ser Pilo Paga raised immediate college enrolment substantially \citep{londono2020}; affects high scorers, not typical NINIs. Others not evaluated. & Medium \\
\midrule
\multicolumn{6}{l}{\textit{Labour-market rules}}\\
Military ID card (\textit{libreta militar}), men only & Law 1738 of 2014: not required for an undergraduate degree. Law 1780 of 2016 art.\ 20: 18 months to define military status once hired; Law 1861 of 2017 art.\ 42: employers may not require it before hiring. & E3 & Men vs.\ women aged 18--24, before vs.\ after 2014/2017 (triple difference by age). & Only a master's thesis found (results not seen). & Medium \\
Youth hiring incentives (ages 18--28) & Law 1780 of 2016 art.\ 7: no family-fund contributions in the first year for new hires aged 18--28 if headcount rises. Law 2155 of 2021 art.\ 24 and Decree 1399 of 2021: subsidy of 25\% of the minimum wage per additional young worker, up to 12 payments, until about 2023 \verify. & E3 & Age discontinuities at 18 and 28/29, before vs.\ after. & No evaluation found. & Medium \\
Minimum wage & +16\% in 2023 and further large increases \verify. & E3 & Minimum-wage bite (Kaitz index) by city and age. & Minimum-wage effects on informality are largest for youth, women and the less educated \citep{banrep1104}. & Medium \\
\midrule
\multicolumn{6}{l}{\textit{Care and transfers}}\\
Early-childhood centres (De Cero a Siempre / CDI) & Rapid expansion of comprehensive-care places 2011--2013; Law 1804 of 2016. & E4 & Place growth by municipality $\times$ young women with vs.\ without children aged 0--5. & ELCA-based evaluation of child outcomes \citep{uniandes2017}; maternal labour only in a master's thesis. & High \\
J\'ovenes en Paz & Law 2294 of 2023 art.\ 348; Decree 1649 of 2023: up to COP 1 million per month for ages 14--28 in prioritised territories; target $\approx$100,000 youth. & E1 & Territorial prioritisation; eligibility rules. & No evaluation found. & Medium \\
Sac\'udete; adolescent work rules; child-marriage ban & Law 2231 of 2022 (framework law, no transfer); Law 1098 of 2006 (work at 15--17 requires a labour-inspector permit); Law 2447 of 2025 (bans marriage and unions under 18). & E3, E4 & Weak or too recent. & --- & Low \\
\end{longtable}
\endgroup

\paragraph{Descriptive evidence to collect first.} Because national reforms lack a clean control group in a repeated cross-section, the first pass should use \emph{eligible vs.\ ineligible} comparisons within the GEIH:
\begin{enumerate}
\item NINI rates by single year of age, sex and completed level, plotted in event time around 2010 (repetition cap), 2012 (free schooling, J\'ovenes en Acci\'on redesign), 2014--2017 (military card rules), 2021 (youth hiring subsidy), 2023 (minimum wage, free tuition) and 2025 (apprenticeship reform).
\item Men vs.\ women aged 18--24 around the military-card changes; ages 17 vs.\ 18 and 28 vs.\ 29 around the youth hiring incentives.
\item Young women with vs.\ without children aged 0--5 in departments with fast vs.\ slow growth of early-childhood places.
\item With administrative data, school-level adoption of SENA articulation (SIMAT) against graduates' later SENA/tertiary enrolment (SPADIES) and formal work (PILA).
\end{enumerate}
A visible break in the eligible group but not in the comparison group is the signal that a policy deserves a full causal design.

% ==================================================================
\section{Literature review}\label{sec:lit}
% ==================================================================

\subsection{Colombia: academic and working papers}

The Colombian literature agrees on the regularities above but none of it is causal. It is methodologically homogeneous: cross-sectional household surveys (GEIH or its predecessors) and multinomial logit models over four states (studies only, works only, both, neither). Only two studies separate home-care NINIs from job-seeking NINIs \citep{hernandez2016,lenis2021}, and only one uses time-use data \citep{ochoa2015}. Our search found no article on Colombian NINIs in \textit{Desarrollo y Sociedad}, \textit{Lecturas de Econom\'ia}, \textit{Ensayos sobre Pol\'itica Econ\'omica} or \textit{Coyuntura Econ\'omica}, nor in the CEDE or Fedesarrollo series (search within journals was not exhaustive).

\begingroup\small
\begin{longtable}{P{3.0cm}P{3.6cm}P{2.6cm}P{5.6cm}}
\caption{Colombian studies on NINIs}\label{tab:litcol}\\
\toprule
Study & Data and period & Method & Main findings \\
\midrule
\endfirsthead
\toprule
Study & Data and period & Method & Main findings \\
\midrule
\endhead
\bottomrule
\endfoot
\citet{lasso2026} & Spliced DANE household surveys (ENH, ECH, GEIH), 1984--2024, seven metro areas, ages 12--40 & Synthetic cohorts; age-period-cohort decomposition; aggregate first differences; cohort-level ``scarring'' regression (correlational) & 1 in 5 is NINI; rate jumps at 16--18; cohort share stays high; countercyclical; 5\% vs.\ 40\% by income quintile; cohort NINI share at 16--22 associated with lower employment (men) and less salaried work (women) at 36--42. Most relevant paper; see Section~\ref{sec:lasso}. \\
\citet{sanchez2021} & DANE household survey (year \verify), ages 14--28 & Multinomial logit & 2.7 M NINIs; gender bias against women; parental education and lower household unemployment reduce the probability. \\
\citet{lopez2020} & Household surveys 2015--2018 & Multinomial logit & 25\% of women vs.\ 8\% of men; children under 5 and being partnered matter most. Most direct on the care channel. \\
\citet{hernandez2016} & GEIH Q3 2012, ages 15--19 & Multinomial logit by sex, three NINI types & Strong gender differentiation; separates job-seeking, home-care and other NINIs. \\
\citet{ochoa2015} & ENUT 2012--13, ages 15--24 & Descriptive (time use) & Female NINIs spend more time on unpaid work; male NINIs on job search. Only study using time-use data. \\
\citet{mora2017} & Cali, 2013 & Unemployment duration models & Being Afro-Colombian, female or living in certain comunas lengthens unemployment. Only within-city ethnic/territorial study. \\
\citet{lenis2021} & GEIH, 13 metro areas, 2014--2019, ages 18--28 & Logit (``panel'' \verify) & Marital status and relationship to household head matter more than income or education; job-seeking and home-care NINIs behave as distinct groups. \\
\citet{ham2021} & Household surveys 2008--2017 & Descriptive & Youth labour-market gaps by gender, region, education; critical review of youth programmes. \\
\citet{castillo2020} & DANE household surveys & Age-period-cohort, probit, multinomial logit & High NINI incidence; education alone does not ensure good labour insertion. \\
\citet{otero2026} & Cartagena, 2012--2024 & Descriptive & Young women have higher unemployment and inactivity; exclusion linked to illegal-activity risk. \\
\citet{henao2024} & Review of 85 documents & Qualitative & Only article focused on female NINIs; claims about causes lack an identification strategy. \\
\citet{bernal2019} & Bogot\'a, ages 15--24 & Descriptive & 2 in 3 NINIs are women; 1 in 5 women aged 15--24 is already a mother. \\
\end{longtable}
\endgroup

\subsection{A closer look at \citet{lasso2026}}\label{sec:lasso}

Because it is the most complete Colombian study, and the natural benchmark for this project, we reviewed it in full.

\paragraph{What it does.} It builds synthetic cohorts from DANE household surveys for the seven main metropolitan areas, 1984--2024, splicing four survey designs (ENH 1984--1999, ECH 2000--2006, GEIH 2006--2020, new GEIH 2021--) by reweighting to population back-projections. It (i) tracks cohort trajectories in four and eight activity categories, including a strict ``Ninini'' category (not studying, not working, not searching), by sex and by income quintile; (ii) decomposes each category into age, cohort and period effects following \citet{deaton2019}, with period effects constrained to be orthogonal to a linear trend; (iii) regresses changes in cohort NINI rates across age groups on aggregate community and macroeconomic variables (36 observations per regression); and (iv) relates labour outcomes at ages 36--42 to the cohort's NINI share 20 years earlier, following \citet{cruces2012} (147 cohort-age-year cells).

\paragraph{Main results.} The age effect for NINI status rises from about 5\% to 22\% between 16 and 18; the cohort effect for NINI is stable since the early-1970s cohorts because rising ``only studies'' is offset by falling ``only works''; the period effect is countercyclical; the female NINI rate stays above 20\% after age 17 but falls across cohorts; and the bottom-quintile rate is about 40\% versus 5\% in the top quintile.

\paragraph{What it says about \emph{why}.} Little that is tested. (a) The main summary, ``low tertiary coverage and low labour-market insertion'', restates the definition of NINI rather than explaining it. (b) Mechanisms for women (early marriage, teen pregnancy, care) and men (dropout followed by informal work) are taken from other studies; the authors note they lack information on unpaid work. (c) Their own interpretations (liquidity constraints of poor families, added-worker behaviour in recessions, possible preference changes in Generation Z, Venezuelan migration behind the stagnation of attendance after 2017) are stated tentatively and left to ``further research''. (d) Section~6 explores correlations with national time series (GDP, inflation, life expectancy, infant mortality, inequality), which cannot discriminate among mechanisms. The paper contains no individual-level analysis.

\paragraph{Limitations relevant for this project.}
\begin{enumerate}
\item \textbf{Stock persistence $\neq$ individual persistence.} Synthetic cohorts show that the cohort NINI share stays high, not that the same people remain NINI.
\item \textbf{Ecological scarring.} The scarring regression relates cohort-level aggregates; with cohort and time effects, the identifying variation is the cohort's own NINI age profile between 16 and 22 mapped onto outcomes at 36--42, partly confounded with the conditions (e.g.\ the 1998--99 recession) under which a cohort was young \likely. It also controls for likely mediators (schooling, partnership, headship, young children).
\item \textbf{Aggregate determinants.} Thirty-six observations, nine time-series regressors and some implausible signs (life expectancy raising NINI) \likely.
\item \textbf{Care not separated}, although the surveys record ``household chores'' as main activity \verify.
\item \textbf{Seven cities only}: 17.2\% in 2024 versus 23\% nationally in OECD data (the paper's own footnote); Caribbean and intermediate cities and rural areas are excluded.
\item \textbf{Income quintiles are endogenous}: current household per-capita income moves with the youth's own work and household formation \likely.
\item \textbf{APC normalisation and inference}: conclusions depend on the Deaton normalisation; clustering on seven cities is fragile; survey redesigns (2000, 2006, 2021; working age moving from 12+ to 15+) and the halved 2020 sample may load onto period or cohort effects.
\end{enumerate}

\paragraph{How this project builds on it.} We (i) replicate their seven-city series as a validation of our pipeline; (ii) extend it nationally and to 23 cities; (iii) decompose NINIs into the typology of Table~\ref{tab:typology}; (iv) estimate individual-level models; (v) measure individual transitions with linked administrative data; and (vi) use policy and local-demand variation to separate the explanations in Section~\ref{sec:why}. A suggestive pattern worth pursuing: higher cohort NINI shares 20 years earlier are associated with \emph{more} non-salaried (informal) work among women (their Table~A.6), suggesting that female scarring may operate through job quality rather than employment.

\subsection{Official and institutional reports}

\begin{itemize}
\item \textbf{DANE} has no thematic report on NINIs. It publishes the monthly bulletin \textit{Mercado laboral de la juventud} (rolling quarters, ages 15--28, broad definition) \citep{dane2026}, a 2021 statistical note on youth \citep{dane2021} and a 2020 sociodemographic overview based on the 2018 census \citep{dane2020}. No ENUT product on NINIs was found.
\item \textbf{DNP} has no study of its own; CONPES 173 and 4040 use DANE diagnostics \citep{conpes173,conpes4040}.
\item \textbf{ICBF} characterised NINIs aged 14--28 into three profiles with policy proposals \citep{icbf}.
\item \textbf{Banco de la Rep\'ublica} is the most analytical institutional source \citep{lasso2026,otero2026}.
\item \textbf{LEE (Javeriana)} published two descriptive reports; the 2024 report uses a strict definition and analyses departments and the transition to tertiary education \citep{lee60,lee99}.
\item \textbf{Local observatories}: Bogot\'a \citep{ode2024}, Cartagena, Medell\'in, Tolima and Manizales publish descriptive notes.
\item \textbf{International}: OECD \citep{oecd2022,oecdgender}, ILO \citep{ilo2024col}, and the IDB survey \citep{novella2018} provide Colombia figures.
\end{itemize}

\subsection{Latin American and international reference literature}

\paragraph{Regional canon.} \citet{dehoyos2016} is the standard reference: about one in five Latin Americans aged 15--24 is NINI, 60\% come from poor or vulnerable households and two thirds are women, yet the increase of about two million NINIs over two decades was driven by men. \citet{cardenas2015} document persistence during the 2000s boom across 18 countries. \citet{behrman2014} provide the dynamic human-capital framework used by most of this literature, and \citet{szekely2015} apply a synthetic-cohort approach across the region, the template followed by \citet{lasso2026}. \citet{novella2018} survey 15,000 youth in nine countries including Colombia. \citet{tornarolli2017} provides the longest harmonised series. The ILO's \textit{Global Employment Trends for Youth 2024} reports a global NEET rate of 20.4\% (28.1\% for women, 13.1\% for men) \citep{iloget2024}. All of this work is descriptive.

\paragraph{Measurement.} \citet{elder2015} argues that the NEET concept is easily misread and proposes excluding the unemployed (NLEET). \citet{eurofound2012} estimate determinants and costs in Europe. \citet{freire2021} estimate separate models for inactive and unemployed \textit{nem-nem} youth in Brazil; separating the two changes the determinants. This is the most direct methodological template for Colombia.

\paragraph{Gender.} \citet{liu2023} find that at least one in four women aged 20--25 is NINI in 12 Latin American countries (census data), with early family formation as the main correlate. The only causal evidence on the marriage channel comes from MENA: \citet{assaad2017} use instrumental variables and find that early marriage reduces women's employment by 16--47\%.

\paragraph{Adjacent causal evidence (none uses NINI status as an outcome).}
\begin{itemize}
\item \textbf{Training:} J\'ovenes en Acci\'on was randomised; women offered the programme earned about 20\% more and were more likely to hold paid jobs \citep{attanasio2011}; formal-sector effects persisted up to eight years \citep{attanasio2017}; trainees and their relatives increased educational attainment \citep{kugler2022}. Meta-analyses find small average effects, larger for women, disadvantaged youth and middle-income countries \citep{kluve2019,card2018}; see also \citet{ibarraran2019} for the Dominican Republic.
\item \textbf{Cash transfers:} Familias en Acci\'on reduced teenage pregnancy, male arrests and male dropout in Medell\'in \citep{attanasio2021} and raised secondary completion \citep{baez2011}.
\item \textbf{Teen motherhood:} \citet{aguero} instruments teen pregnancy with age at menarche in Colombia.
\item \textbf{Childcare:} a Chilean experiment raises maternal employment \citep{martinez2017}.
\item \textbf{Scarring:} early unemployment has lasting effects on later employment and wages \citep{gregg2001,mroz2006}; NEET spells predict worse occupational outcomes over 20 years \citep{ralston2016}. This literature concerns mostly male unemployment in rich countries, not female inactivity due to care.
\item \textbf{Crime:} in Mexico, the association between idle young men and homicides appears only after 2007 \citep{dehoyos2016crime} (state panel, not causal).
\end{itemize}

\subsection{Gaps}

\begin{enumerate}
\item No causal evidence on the determinants of entering NINI status in Colombia.
\item No exogenous variation used to study motherhood, partnership and care.
\item No individual-level evidence on scarring with administrative data, and none on female care-related inactivity anywhere in middle-income countries.
\item No evaluation of current policies (Sac\'udete, J\'ovenes en Paz, free tuition, Renta Joven) with NINI status as an outcome.
\item No study of Venezuelan migrant youth.
\item Ethnicity and region covered only for Cali and Cartagena.
\item No academic decomposition of why Colombia's female rate is so high relative to the OECD.
\item Existing randomised evaluations (J\'ovenes en Acci\'on, Familias en Acci\'on) have never been re-analysed with disaggregated NINI status as an outcome; this is the cheapest contribution, since the variation already exists \likely.
\end{enumerate}

% ==================================================================
\section{Data}\label{sec:data}
% ==================================================================

\begingroup\small
\begin{longtable}{P{3.2cm}P{7.2cm}P{3.6cm}}
\caption{Data sources}\label{tab:data}\\
\toprule
Source & What it provides & Access \\
\midrule
\endfirsthead
\toprule
Source & What it provides & Access \\
\midrule
\endhead
\bottomrule
\endfoot
GEIH (DANE) & Monthly labour-force survey: attendance, labour status, main activity, reasons for inactivity, household roster, migration module, region. Basis for the typology. Use back-cast series after the 2022 redesign. & Public microdata \\
ENUT (DANE) & Time use (2012--13, 2016--17, 2020--21): hours of unpaid care; validates the care group. & Public microdata \\
ECV (DANE) & Annual living-conditions survey: education, household composition. & Public microdata \\
ENDS & Fertility and partnership histories; age at first union and first birth. & Public (2015; check 2025 round) \\
ELCA (Uniandes) & Household panel following youth over time. & Public with registration \verify \\
Census 2018 & Attendance and activity for the full population; municipal rates by ethnicity. & Public microdata \\
ETET (ILO--DANE 2013, 2015) & School-to-work transition survey; almost unused for NINIs. & Public microdata \\
Saber 11 (ICFES) & Cohorts of secondary graduates; scores and socioeconomic background. & Public; linkage restricted \\
SPADIES / SNIES, SENA & Tertiary and technical enrolment and dropout. & Restricted \\
PILA & Formal employment and wages (social-security contributions). & Restricted (DANE secure room / agreements) \\
SIMAT & School enrolment. & Restricted \\
SISBEN IV & Targeting score and programme eligibility cut-offs. & Restricted \\
Vital statistics (births) & Mother's age and characteristics; early fertility. & Public microdata \\
\end{longtable}
\endgroup

% ==================================================================
\section{Proof of concept: a first empirical step}\label{sec:plan}
% ==================================================================

\subsection{Stage 1 -- Descriptive decomposition (weeks)}

Using GEIH 2022--2025 (redesigned survey) plus back-cast series:
\begin{enumerate}
\item Replicate DANE's broad rate for ages 15--28 by quarter.
\item Build the typology in Table~\ref{tab:typology} and report the stock by group, sex, single year of age and presence of own children under five.
\item Produce a sensitivity table: broad vs.\ strict definition; ages 14--28, 15--28, 15--24 and 18--24; national vs.\ 13 cities.
\item Validate the care group with ENUT hours.
\item Plot the age profile 14--28 by sex and completed level, to test whether the 16--18 jump coincides with graduation age (Q2).
\item Replicate the seven-city series of \citet{lasso2026} and extend it nationally and to 23 cities.
\item Estimate an individual-level multinomial model of NINI type (Table~\ref{tab:typology}) on personal and household characteristics, using predetermined socioeconomic proxies (household head's education, estrato) instead of current income.
\item Produce the event-time and eligible-vs-ineligible plots listed at the end of Section~\ref{sec:policy}.
\end{enumerate}

\noindent\textbf{What this stage answers.} Whether the gender paradox is mostly care (Q1); whether the 16--18 jump is a graduation-age artefact (Q2); how much of the gap with the OECD is purely definitional (Q5); and which policy changes in Table~\ref{tab:policy} leave a visible break in the eligible group. The answers determine which explanation in Table~\ref{tab:why} deserves a causal design.

\subsection{Stage 2 -- Candidate causal designs (to be selected after Stage 1)}

\begin{table}[htbp]
\centering\small
\caption{Candidate designs}\label{tab:designs}
\begin{tabularx}{\linewidth}{P{3.6cm}LP{3.2cm}}
\toprule
Design & Idea & Main risk \\
\midrule
Re-analysis of existing RCTs & Use J\'ovenes en Acci\'on / Familias en Acci\'on randomisation with disaggregated NINI status as outcome & Access to original data; outcome availability \\
SISBEN RD of the 2012 J\'ovenes en Acci\'on / Renta Joven transfer & Extend \citet{gamboa2021} from SENA enrolment to NINI status and formal work (PILA) & Access to SISBEN scores linked to outcomes \\
Repetition-cap repeal (2010) & Extend \citet{ferreira2020}: follow induced dropouts into work vs.\ NINI (SIMAT--PILA) & Linkage; school-level confounders \\
Staggered SENA articulation & Staggered DiD by school adoption on graduates' NINI status & Endogenous adoption; coverage data \\
Military-card rules & Triple difference: men vs.\ women, ages 18--24 vs.\ older, before vs.\ after 2014/2017 & Concurrent shocks to young men \\
Eligibility cut-offs & RD around SISBEN thresholds of youth programmes; outcome: NINI status in GEIH-linked or administrative data & Linking eligibility to outcomes \\
Childcare supply shocks & Staggered expansion of early-childhood centres (CDIs) by municipality as a shock to young mothers' care constraint & Endogenous placement \\
Graduation-age discontinuity & Compare youths born just before/after school-entry cut-offs, who graduate before or after 18 & Cut-off enforcement; data on birth month \\
Scarring via cohort conditions & Local labour demand at the year of leaving school as an instrument for NINI spells, outcomes in PILA & Exclusion restriction \\
\bottomrule
\end{tabularx}
\end{table}

\subsection{Risks}

\begin{itemize}
\item Restricted administrative data require formal agreements and time.
\item The GEIH reasons-for-inactivity questions may have changed with the 2022 redesign; the typology must be harmonised across questionnaires.
\item Several figures in this note come from press summaries; Stage 1 will replace them with own estimates.
\end{itemize}

% ==================================================================
\section{Next steps}
% ==================================================================

\begin{enumerate}
\item Read in full: \citet{lasso2026} (tables by sex and income, data, pseudo-panel strategy); \citet{sanchez2021}; \citet{lopez2020}; \citet{lee99}; \citet{icbf}; the latest DANE bulletin and \citet{dane2021}; \citet{aguero}.
\item Locate: the original ILO 2024 report; the technical document of \citet{bernal2019}; the LaboUR (Rosario) report; the 2026 \textit{Cuadernos de Econom\'ia} article on NINIs aged 14--28 (authors not identified).
\item Search directly in Google Scholar, RePEc, Documentos CEDE, Fedesarrollo and university thesis repositories.
\item Download GEIH 2022--2025 and ENUT microdata and run Stage 1.
\end{enumerate}

\appendix
% ==================================================================
\section{Sketch of the typology in Stata}
% ==================================================================

The variable names below follow the pre-2022 GEIH questionnaire and are indicative only; they must be checked against the dictionary of the redesigned GEIH.

\begin{lstlisting}[language={}]
* Indicative sketch -- verify variable names in the current GEIH dictionary
* P6040 age | P6020 sex | P6170 attends school | P6240 main activity last week
* OCI employed | DSI unemployed | INI inactive | reason for not seeking: see questionnaire

keep if inrange(P6040, 15, 28)
gen nini_broad  = (P6170 == 2) & (OCI != 1)
gen nini_strict = nini_broad & (INI == 1)

gen     nini_type = .
replace nini_type = 1 if nini_broad & DSI == 1                          // searching
replace nini_type = 2 if nini_strict & P6040 <= 19 & completed_secondary // transition
replace nini_type = 3 if nini_strict & P6240 == 4 & missing(nini_type)  // household / care
replace nini_type = 4 if nini_strict & P6240 == 5 & missing(nini_type)  // disability
replace nini_type = 5 if nini_strict & discouraged & missing(nini_type) // discouraged
replace nini_type = 6 if nini_broad & missing(nini_type)               // disconnected

label define ntype 1 "Searching" 2 "Transition" 3 "Care" 4 "Disability" ///
                   5 "Discouraged" 6 "Disconnected"
label values nini_type ntype

* Shares by sex and single year of age, using expansion factors
tab nini_type P6020 [iw = fex_c] if nini_broad, col
\end{lstlisting}

\newpage
\bibliographystyle{plainnat}
\bibliography{ninis_refs}

\end{document}
